% Einheit 4 von 5 – Gleichung bestimmen (bank/lineare-funktionen/e4.jsonl, zusammenbau v0.1)
%% TODO Zweigzeile Teil 1: was hier gelernt wird (Satz in Schülersprache aus dem Einheitstitel) – Einheit 4
\einheitenkopf[e4]{Einheit 4 von 5 · Gleichung bestimmen}
\zweigzeile{ab Kl. 8 · P10}
\setcounter{aufgabe}{23}

%% TODO Titel: Ich-kann-Satz fehlt in der Bank (Platzhalter: Kettenname) – Nr. 24
%% TODO Anweisung: ein Satz über der ersten Teilaufgabe fehlt in der Bank – Nr. 24
\begin{aufgabe}{Gleichung bestimmen}
\swz{Gleichung? Bestimme die Gleichung der Geraden g aus dem Graphen.}{\begin{ksys}[xmin=-5,xmax=5,ymin=-5,ymax=5,ablesen] \gerade{2}{-3}{g} \end{ksys}}
\swz{Gleichung? Bestimme die Gleichung der Geraden g aus dem Graphen.}{\begin{ksys}[xmin=-5,xmax=5,ymin=-5,ymax=5,ablesen] \gerade{-1}{4}{g} \end{ksys}}
\swz{Gleichung? Bestimme die Gleichung der Geraden g aus dem Graphen.}{\begin{ksys}[xmin=-5,xmax=5,ymin=-5,ymax=5,ablesen] \gerade{3}{1}{g} \end{ksys}}
\swz{Gleichung? Bestimme die Gleichung der Geraden g aus dem Graphen.}{\begin{ksys}[xmin=-5,xmax=5,ymin=-5,ymax=5,ablesen] \gerade{-2}{-1}{g} \end{ksys}}
\swz{Gleichung? Bestimme die Gleichung der Geraden g aus dem Graphen.}{\begin{ksys}[xmin=-5,xmax=5,ymin=-5,ymax=5,ablesen] \gerade{1}{-4}{g} \end{ksys}}
\swfrage{Was bleibt gleich, was ändert sich?}
%% TODO Darstellung für schwach fehlt in der Bank (lineare-funktionen-e4-k1-s2-v1; Abschnitt „Für schwache Schüler“)
\swz{$\text{Gleichung? Eine Gerade hat die Steigung $m = 2$ und geht durch $P(4|7)$. Bestimme ihre Gleichung.}$}{}
%% TODO Darstellung für schwach fehlt in der Bank (lineare-funktionen-e4-k1-s3-v1; Abschnitt „Für schwache Schüler“)
\swz{$\text{Gleichung aus zwei Punkten: Bestimme die Gleichung der Geraden durch $A(0|4)$ und $B(2|10)$.}$}{}
%% TODO Darstellung für schwach fehlt in der Bank (lineare-funktionen-e4-k1-s4-v1; Abschnitt „Für schwache Schüler“)
\swz{$\text{Gleichung aus zwei Punkten: Bestimme die Gleichung der Geraden durch $A(2|3)$ und $B(5|9)$.}$}{}
%% TODO Darstellung für schwach fehlt in der Bank (lineare-funktionen-e4-k1-s5-v1; Abschnitt „Für schwache Schüler“)
\swz{$\text{Gleichung aus zwei Punkten: Bestimme die Gleichung der Geraden durch $A(1|6)$ und $B(3|0)$.}$}{}
\swa{Gerade durch A und B: Zeichne die Gerade f durch $A(-2|8)$ und $B(4|-1)$, entscheide für jede Aussage, ob sie wahr oder falsch ist, und gib eine Gleichung von f an. \hfill (P10 2025 OS) \\ \kreuz{wahr}\kreuz{falsch} f verläuft monoton steigend. \\ \kreuz{wahr}\kreuz{falsch} f schneidet die y-Achse in $(0|5)$. f(x) = \leerfeld}{\begin{ksys}[xmin=-4,xmax=5,ymin=-4,ymax=9] \end{ksys}}
\end{aufgabe}

%% TODO Titel: Ich-kann-Satz fehlt in der Bank (Platzhalter: Kettenname) – Nr. 25
%% TODO Anweisung: ein Satz über der ersten Teilaufgabe fehlt in der Bank – Nr. 25
\begin{aufgabe}{Gerade durch zwei Punkte zeichnen}
\swa{Gerade durch zwei Punkte: Trage $A(-3|-1)$ und $B(2|4)$ ein und zeichne die Gerade durch beide Punkte.}{\begin{ksys}[xmin=-4,xmax=4,ymin=-4,ymax=5] \end{ksys}}
\end{aufgabe}

%% TODO Titel: Ich-kann-Satz fehlt in der Bank (Platzhalter: Kettenname) – Nr. 26
%% TODO Anweisung: ein Satz über der ersten Teilaufgabe fehlt in der Bank – Nr. 26
\begin{aufgabe}{Schnittpunkt zweier Geraden rechnerisch}
%% TODO Darstellung für schwach fehlt in der Bank (lineare-funktionen-e4-k3-s1-v1; Abschnitt „Für schwache Schüler“)
\swz{$\text{Schnittpunkt? Berechne den Schnittpunkt der Geraden $f(x) = 3x - 4$ und $g(x) = -x + 8$.}$}{}
\end{aufgabe}

%% TODO Titel: Ich-kann-Satz fehlt in der Bank (Platzhalter: Kettenname) – Nr. 27
%% TODO Anweisung: ein Satz über der ersten Teilaufgabe fehlt in der Bank – Nr. 27
\begin{aufgabe}{Gleichung bestimmen – Fehler finden, Begründen, Darstellungswechsel, Anwendung}
\swz[3]{Ich finde den Fehler: Für $A(2|3)$ und $B(5|12)$ rechnet Kim: \rechnung{m &= (3 - 12) : (5 - 2) \\ &= -3} Benenne den Fehler und rechne richtig.}{}
\swfrage{Schneiden sich die Geraden $f(x) = 3x - 7$ und $g(x) = 3x + 5$? Begründe ohne Rechnung.}
\swz{Gleichung zur Geraden: Bestimme die Gleichung der Geraden g aus dem Graphen.}{\begin{ksys}[xmin=-5,xmax=5,ymin=-5,ymax=5,ablesen] \gerade{-0.5}{1}{g} \end{ksys}}
\swz[3]{$\text{Wie hoch ist die Grundgebühr? Für eine Taxifahrt von 4 km zahlt man 12{,}60 €, für 10 km 25{,}20 €. Der Preis setzt sich aus Grundgebühr und festem Preis je km zusammen.}$}{}
\end{aufgabe}

\uebersichtskasten{Steigung aus zwei Punkten: m = (y₂ − y₁) : (x₂ − x₁) \quad A(1 | 3), B(3 | 7): m = (7 − 3) : (3 − 1) = 2 \\ Dann n: m und einen Punkt in y = m · x + n einsetzen, nach n auflösen. \\ Schnittpunkt zweier Geraden: Terme gleichsetzen, x ausrechnen, y aus einer Gleichung.}
